\documentclass[10pt]{article}

% Specify the margins and text size.
\setlength{\textwidth}{6.4in}
\setlength{\textheight}{9.5in}
\setlength{\oddsidemargin}{0pt}
\setlength{\evensidemargin}{0pt}
\setlength{\topmargin}{0pt}
\setlength{\hoffset}{.05in}
\setlength{\voffset}{-.8in}

\setlength{\parskip}{5pt}
\setlength{\parindent}{0pt}

% Load some fonts and symbol packages
\usepackage{latexsym}
\usepackage{pifont}       % contains 'star' symbol for counterinsurgency handbook title
\usepackage{yfonts} 
\usepackage{amsmath}
\usepackage{amsfonts}

\usepackage{graphicx}     % actually, this is loaded by pstricks
\usepackage[T1]{fontenc}
\usepackage{ifthen}
\usepackage{pstricks,pst-func,pst-text,pst-node,multido,pst-plot,calc,pst-3dplot}
%\usepackage[all]{xy}
%\usepackage{animate}

% The hyperref package inserts links into the pdf.
\definecolor{MyLinkColor}{rgb}{.1,.2,1}
\definecolor{MyCiteColor}{rgb}{.1,1,.2}
\definecolor{MyURLColor}{rgb}{.4,.4,.4}
\usepackage[backref=true,pagebackref=false,hyperindex,colorlinks=true,
  linkcolor=MyLinkColor,urlcolor=MyURLColor]{hyperref}


% The tweaklist package is something I found on the web.  It provides a simple interface
% for making changes to spacing used in the itemize and enumerate environments.  Comment
% this out if you don't care to use tweaklist.
\usepackage{tweaklist}
\renewcommand{\itemhook}{\setlength{\parskip}{2pt}\setlength{\parsep}%
{1pt}\setlength{\topsep}{0pt}\setlength{\itemsep}{0pt}}

\newcommand{\N}{\ensuremath\mathbb{N}}
\newcommand{\Q}{\ensuremath\mathbb{Q}}
\newcommand{\R}{\ensuremath\mathbb{R}}
\newcommand{\C}{\ensuremath\mathbb{C}}
\newcommand{\z}{\hphantom{0}}

\begin{document}
\pagestyle{empty}

\href{http://www.su.edu}{\includegraphics[height=1.75cm]{sulogo.eps}}
\vspace{-1.69cm}

{\small\hfill
\begin{tabular}{cl}
 & Math 207\\ & Introduction to Statistics\\
\end{tabular}
}

\setlength{\baselineskip}{1.05\baselineskip}

\begin{center}
\textbf{\large Quiz:  Mean, Quantiles and Standard Deviation}
\end{center}

1. For each list of numbers, calculate the mean, median and sample
standard deviation.

\hspace{10pt} a) 9, 9, 10, 10, 10, 12
\vspace{.45in}

\hspace{10pt} b) 7, 8, 10, 11, 11, 13
\vspace{.45in}

2. For each list of numbers, calculate the mean, median and sample
standard deviation.  

\hspace{10pt} a) 1, 3, 4, 5, 7
\vspace{.45in}

\hspace{10pt} b) 6, 8, 9, 10, 12
\vspace{.45in}

\hspace{10pt} c)  How are the lists in this exercise related to each other?
How does that property carry over to the mean, median and sd?
\vspace{.7in}


3. Repeat exercise 2 for the following two lists.

\hspace{10pt} a) 1, 3, 4, 5, 7
\vspace{.6in}

\hspace{10pt} b) 3, 9, 12, 15, 21
\vspace{.6in}

\vfill
\eject

\end{document}
